\documentclass[10pt,a4paper]{amsart}
\usepackage{amsmath,amssymb,amsthm,graphicx}
\usepackage{verbatim}
\newenvironment{reference}{\comment}{\endcomment}
%\newenvironment{reference}{}{}
\newenvironment{slogan}{\comment}{\endcomment}
\newenvironment{history}{\comment}{\endcomment}

% For commutative diagrams we use Xy-pic
\usepackage[all]{xy}

% We use 2cell for 2-commutative diagrams.
\xyoption{2cell}
% Theorem environments.
%
\theoremstyle{plain}
\newtheorem{theorem}[subsection]{Theorem}
\newtheorem{proposition}[subsection]{Proposition}
\newtheorem{lemma}[subsection]{Lemma}

\theoremstyle{definition}
\newtheorem{definition}[subsection]{Definition}
\newtheorem{example}[subsection]{Example}
\newtheorem{exercise}[subsection]{Exercise}
\newtheorem{situation}[subsection]{Situation}

\theoremstyle{remark}
\newtheorem{remark}[subsection]{Remark}
\newtheorem{remarks}[subsection]{Remarks}

\numberwithin{equation}{subsection}

% Macros
%
\def\lim{\mathop{\mathrm{lim}}\nolimits}
\def\colim{\mathop{\mathrm{colim}}\nolimits}
\def\Spec{\mathop{\mathrm{Spec}}}
\def\Hom{\mathop{\mathrm{Hom}}\nolimits}
\def\Ext{\mathop{\mathrm{Ext}}\nolimits}
\def\SheafHom{\mathop{\mathcal{H}\!\mathit{om}}\nolimits}
\def\SheafExt{\mathop{\mathcal{E}\!\mathit{xt}}\nolimits}
\def\Sch{\mathit{Sch}}
\def\Mor{\mathop{\mathrm{Mor}}\nolimits}
\def\Ob{\mathop{\mathrm{Ob}}\nolimits}
\def\Sh{\mathop{\mathit{Sh}}\nolimits}
\def\NL{\mathop{N\!L}\nolimits}
\def\CH{\mathop{\mathrm{CH}}\nolimits}
\def\proetale{{pro\text{-}\acute{e}tale}}
\def\etale{{\acute{e}tale}}
\def\QCoh{\mathit{QCoh}}
\def\Ker{\mathop{\mathrm{Ker}}}
\def\Im{\mathop{\mathrm{Im}}}
\def\Coker{\mathop{\mathrm{Coker}}}
\def\Coim{\mathop{\mathrm{Coim}}}

% Boxtimes
%
\DeclareMathSymbol{\boxtimes}{\mathbin}{AMSa}{"02}

%
% Macros for moduli stacks/spaces
%
\def\QCohstack{\mathcal{QC}\!\mathit{oh}}
\def\Cohstack{\mathcal{C}\!\mathit{oh}}
\def\Spacesstack{\mathcal{S}\!\mathit{paces}}
\def\Quotfunctor{\mathrm{Quot}}
\def\Hilbfunctor{\mathrm{Hilb}}
\def\Curvesstack{\mathcal{C}\!\mathit{urves}}
\def\Polarizedstack{\mathcal{P}\!\mathit{olarized}}
\def\Complexesstack{\mathcal{C}\!\mathit{omplexes}}
% \Pic is the operator that assigns to X its picard group, usage \Pic(X)
% \Picardstack_{X/B} denotes the Picard stack of X over B
% \Picardfunctor_{X/B} denotes the Picard functor of X over B
\def\Pic{\mathop{\mathrm{Pic}}\nolimits}
\def\Picardstack{\mathcal{P}\!\mathit{ic}}
\def\Picardfunctor{\mathrm{Pic}}
\def\Deformationcategory{\mathcal{D}\!\mathit{ef}}
\usepackage{cleveref}
\crefname{originalreference}{original reference}{original references}
\Crefname{originalreference}{Original reference}{Original references}
\makeatletter
\newlabel{lemma-linear-series}{{0CCP}{source}{}{sourcepointer-1}{}}
\newlabel{lemma-linear-series@cref}{{[originalreference][0][]0CCP}{[1][0][]source}{}{}{}}
\makeatother
\crefname{originallemma}{original lemma}{original lemmas}
\Crefname{originallemma}{Original lemma}{Original lemmas}
\makeatletter
\newlabel{varieties-lemma-check-invertible-sheaf-trivial}{{0B40}{source}{}{sourcepointer-2}{}}
\newlabel{varieties-lemma-check-invertible-sheaf-trivial@cref}{{[originallemma][0][]0B40}{[1][0][]source}{}{}{}}
\makeatother
\crefname{originallemma}{original lemma}{original lemmas}
\Crefname{originallemma}{Original lemma}{Original lemmas}
\makeatletter
\newlabel{varieties-lemma-degree-pullback-map-proper-curves}{{0AYZ}{source}{}{sourcepointer-3}{}}
\newlabel{varieties-lemma-degree-pullback-map-proper-curves@cref}{{[originallemma][0][]0AYZ}{[1][0][]source}{}{}{}}
\makeatother
\title{A pencil on a nonsingular proper curve gives either a map of degree at most d or a trivial map with nontrivial global functions}
\author{The Stacks Project Authors}
\date{stacks-project:curves.tex:local-snapshot-20261012; GNU FDL 1.2 or later}
\begin{document}
\maketitle
\noindent\textit{Editorial note.} All delivered mathematical text is the author's own, copied from the pinned Stacks Project chapter file curves.tex; no statement or proof step was rewritten, reordered or removed.\par\medskip
\noindent\textit{Editorial note.} The counted proof is the author's complete argument: the degree bounds, the trivial case with the H\textasciicircum{}0 consequence, and the nonconstant case with the degree formula are all retained verbatim.\par\medskip
\noindent\textit{Editorial note.} The label lemma-g1d names the counted statement, and the proof environment immediately following it in the source is delivered verbatim as its proof.\par\medskip
\noindent\textit{Editorial note.} The counted statement is carried by the author's own lemma environment in the chapter source.\par\medskip
\noindent\textit{Allowed standard prerequisites (editorial).}\begin{itemize}
\item The morphism attached to a pencil. Same chapter, Lemma lemma-linear-series. Produces the morphism phi: X -> P\textasciicircum{}1 of the g\textasciicircum{}1\_d with L' = phi\textasciicircum{}*O(1) receiving a nonzero map from the pencil.
\item Degree bookkeeping for invertible sheaves. Varieties, Lemmas varieties-lemma-check-invertible-sheaf-trivial and varieties-lemma-degree-pullback-map-proper-curves. Bounds 0 <= deg(L') <= d, forces L' trivial at degree zero, and gives the degree formula for the pullback under a nonconstant map.
\end{itemize}\medskip
\ifdefined\hypertarget\hypertarget{sourcepointer-1}{}\fi
\noindent\textit{Source reference.} Original reference 0CCP: sources/stacks-curves/pinned-original-source.tex, source lines 583--583. Prerequisite pointer invoked at proof line(s) 629; it names a standard fact of the same chapter that the counted argument uses, and every step of the counted proof built on it is retained verbatim in the delivered body.\par\medskip
\ifdefined\hypertarget\hypertarget{sourcepointer-2}{}\fi
\noindent\textit{Source reference.} Original lemma 0B40: sources/stacks-chow/tags.txt, source lines 13983--13983. Prerequisite pointer invoked at proof line(s) 642; it names a standard fact of a sibling chapter that the counted argument uses, and every step of the counted proof built on it is retained verbatim in the delivered body.\par\medskip
\ifdefined\hypertarget\hypertarget{sourcepointer-3}{}\fi
\noindent\textit{Source reference.} Original lemma 0AYZ: sources/stacks-chow/tags.txt, source lines 13807--13807. Prerequisite pointer invoked at proof line(s) 654; it names a standard fact of a sibling chapter that the counted argument uses, and every step of the counted proof built on it is retained verbatim in the delivered body.\par\medskip
\begin{lemma}
\label{lemma-g1d}
Let $k$ be a field. Let $X$ be a nonsingular proper curve over $k$.
Let $(\mathcal{L}, V)$ be a $\mathfrak g^1_d$ on $X$. Then the morphism
$\varphi : X \to \mathbf{P}^1_k$ of Lemma \ref{lemma-linear-series}
either
\begin{enumerate}
\item is nonconstant and has degree $\leq d$, or
\item factors through a closed point of $\mathbf{P}^1_k$ and in this
case $H^0(X, \mathcal{O}_X) \not = k$.
\end{enumerate}
\end{lemma}

\begin{proof}
By Lemma \ref{lemma-linear-series} we see that
$\mathcal{L}' = \varphi^*\mathcal{O}_{\mathbf{P}^1_k}(1)$
has a nonzero map $\mathcal{L}' \to \mathcal{L}$.
Hence by Varieties, Lemma \ref{varieties-lemma-check-invertible-sheaf-trivial}
we see that $0 \leq \deg(\mathcal{L}') \leq d$.
If $\deg(\mathcal{L}') = 0$, then the same lemma tells us
$\mathcal{L}' \cong \mathcal{O}_X$ and since we have
two linearly independent sections we find we are in case (2).
If $\deg(\mathcal{L}') > 0$ then $\varphi$ is nonconstant (since the
pullback of an invertible module by a constant morphism is trivial).
Hence
$$
\deg(\mathcal{L}') =
\deg(X/\mathbf{P}^1_k) \deg(\mathcal{O}_{\mathbf{P}^1_k}(1))
$$
by Varieties, Lemma \ref{varieties-lemma-degree-pullback-map-proper-curves}.
This finishes the proof as the degree of
$\mathcal{O}_{\mathbf{P}^1_k}(1)$ is $1$.
\end{proof}

\end{document}
